\documentclass[12pt]{article}
\usepackage[american]{circuitikz}
\usepackage{amsmath}
\usepackage{geometry}
\geometry{margin=2cm}

\begin{document}
\section{ejemplos de circuitos electricos}
\subsection*{circuito 1: Fuente DC y resitencias}

\centering
\begin{circuitikz}
\draw
(0,0) to [battery, l=$5v$] (0,3)
    to[R, l=$220 \omega$] (3,3)
    --(3,0) --(0,0)
 \end{circuitikz}

\subsection*{circuito 2: Resistencia y capacitor en serie}
\begin{circuitikz}
\draw
(0,0) to [sV, l=$v(t)$] (0,3)
    to[R, l=$4 \micro F$] (6,3)
     --(6,0) --(0,0)
 \end{circuitikz}
 
 \subsection*{circuito 3: Inductor y resistencia en serie}
\begin{circuitikz}
\draw
(0,0) to [battery1, l=$5v$] (0,3)
    to[L, l=$5mH$] (3,3)
    to[R, l=$220 \omega$] (6,3)
     --(6,0) --(0,0)
 \end{circuitikz}

 \subsection*{circuito 4: LEDs en paralelo}
\begin{circuitikz}
\draw
(0,0) to [battery1, invert, l=$V_s$] (0,4)
\draw (0,4)--(4,4)
\draw (0,0)--(4,0)
\draw (1,4) to [led, l=R1] (1,0)
\draw (4,4) to [led, l=R2] (4,0)
\end{circuitikz}




\end{document}